% ===== Chapter 5: auto-converted from Word =====

\section{Discussion}

This chapter interprets the empirical findings of the results and situates them within the theoretical and empirical literature reviewed in Chapter 2. Discussing the main finding of the risk-adjusted outperformance across all five ownership quintiles. Showing the constraints of the Fama French factors when controlling for the returns of the high ownership firms and mentioning limitations of this paper. 


\subsection{The Universe Level Alpha}

That every ownership quintile within the Wood (2025) universe earns positive five factor alpha implies that the risk-adjusted outperformance is not bounded to firms with extreme ownership concentration. The finding extends Wood's return evidence, in which owner operated firms compounded at 10.9\% per year against 6.3\% for the S\&P Global 1200 over three decades. Wood did not adjust the return for specific risk factors and left the possibility that the outperformance reflected compensation for systematic risk open. The five factor regressions reject that explanation: even after controlling for market, size, value, profitability, and investment exposures, the universe earns substantial alpha relative to the developed markets factors.

The individual quintile alphas should not be read as a pure marginal effect of ownership concentration. The alpha estimated for each quintile portfolio combines two effects. The first is a universe level premium shared by every firm in the sample by belonging to the Wood (2025) universe. The second is a within universe concentration premium that varies with ownership across quintiles. Both effects are present in every quintile alpha. The 0.92\% monthly alpha of P1 and the 1.45\% alpha of P5 each include the same universe level contribution. The HML spread, constructed as P5 minus P1, cancels the universe level component because it appears equally in both legs and isolates the concentration premium alone. The HML alpha of 0.53\% per month is therefore the appropriate estimate for testing whether ownership concentration is priced within the universe. The individual quintile alphas address the broader question of whether the universe as a whole outperforms and extends Wood's paper by adjusting for risk factors. Both findings hold in the data and should be reported as distinct results.


\subsection{The Concentration Premium and the Factor Mismatch}

The HML spread captures the additional risk-adjusted return associated with moving from low to high insider ownership within the universe. It survives the momentum factor, the split of the sample, the alternative tercile and decile sorts, and the switch from equal to value weighting. The decile sort shows that the relationship is positive across the ownership distribution. The breakdown of monotonicity at the quintile P4 level is limited to that one cut off and does not appear at coarser or finer partitions. 

Interpreting the alphas requires understanding what the five factor regression actually controls for. The estimated alpha is not the excess return over the broad market index. It is the residual return after controlling for the combination of the five Fama French factors. The Fama French factors were constructed with returns in the developed markets universe, which is dominated by dispersed ownership firms. They successfully capture the systematic exposures common in that universe. They are less effective when used on the return of firms selected on a characteristic like, longer investment horizons, governance structure, and concentration in markets where high insider ownership is common. This therefore does not fully reflect the Wood universe.

This consideration is particularly relevant for the larger regional alphas. The North American HML estimate corresponds to roughly 21\% per year. A return of this size is well above what would be expected from any genuine investable anomaly in a mature market and demands cautious framing. One issue is the North American subsample contains only 141 firms, leaving approximately 28 firms in each tail portfolio. The t-statistic of 1.681 reflects the resulting noise. Next to that, the implementation costs that an investor would face are substantial. Many sample firms have low free floats because insiders hold the majority. Bid-ask spreads on mid-cap names are wide, low-ownership names within the universe are often hard to short, and rebalancing generates turnover costs, stamp duties, and currency conversion frictions. The alphas in this thesis should therefore be understood as gross statistical premia, not as estimates of investable returns net of transaction costs and capacity constraints. They are upper bounds on what a real portfolio strategy could earn. They combine an alignment driven return component with compensation for systematic dimensions and the noise that comes with characteristic based portfolio sorts.


\subsection{Theoretical Mechanisms}

The pattern of results provides a basis for evaluating the competing theoretical perspectives developed in Chapter 2. The convergence of interest hypothesis of Jensen and Meckling (1976) predicts a positive monotonic relationship between insider ownership and risk-adjusted returns. The entrenchment hypothesis associated with Fama and Jensen (1983) and Morck, Shleifer, and Vishny (1988) predicts a non-monotonic relationship. They show that alignment effects are overwhelmed at the highest ownership levels by reduced external discipline. Stewardship theory predicts a universe level premium consistent with insider control fostering long term orientation that benefits shareholders broadly (Donaldson and Davis, 1991).

The evidence is most consistent with a combination of convergence of interest and stewardship perspectives. The grading of alpha across the ownership distribution, supports the convergence prediction. The substantial alpha earned by every quintile supports the broader stewardship interpretation. The entrenchment hypothesis receives no support: the highest ownership decile, containing firms with insider stakes between 68\% and 99.8\%, earns the largest alpha in the analysis. There is no evidence within this sample of a turning point at which entrenchment dominates alignment. These findings extend the existing return based ownership literature to a globally diversified setting. Fahlenbrach (2009) documents founder-CEO outperformance and Lilienfeld-Toal and Ruenzi (2014) document a CEO ownership premium, both in the United States. Anderson and Reeb (2003) and Andres (2008) document family firm outperformance in the United States and continental Europe respectively. This paper generalises to aggregate insider ownership and to a globally diversified universe.


\subsection{The Regional Contrast}

The divergence between the North American and European HML alphas can be understood through the cross country institutional framework. La Porta, Lopez de Silanes, and Shleifer (1999) and Claessens, Djankov, and Lang (2000) document that concentrated insider ownership is the modal structure outside the United States and the United Kingdom. In North America, high insider ownership is rare: a firm with 50\% insider equity shows strong information about governance and managerial commitment. In continental Europe, concentrated ownership is common and less differentiated. The marginal informational content of high ownership is therefore lower in Europe. The North American concentration premium within in the Wood universe is correspondingly compressed. As discussed in Section 5.2, the magnitude of the North American HML estimate must also be read against the small subsample size and the implementation frictions discussed there.

The absence of a significant European HML alpha does not imply European high ownership firms underperform. The European quintile alphas are positive and significant. What is absent is the additional concentration premium seen in North America. This pattern is consistent with Andres (2008) and Barontini and Caprio (2006), who find positive performance effects of concentrated ownership in continental Europe at smaller outcomes than the corresponding United States studies. The weak HML alphas in Asia-Pacific and emerging markets ex-Asia can also be explained by the previous argument: in regions where concentrated ownership is common, the marginal informational content of ownership variation is reduced.


\subsection{Limitations}

Several limitations have to be taken into account in the interpretation of these results. First, the universe construction selects firms on a 20\% ownership threshold. The analysis cannot directly compare high ownership firms to a control group of low ownership firms drawn from the same global universe. It documents how the return premium grades with concentration within the high ownership universe but cannot isolate the effect of moving from zero to higher insider ownership.

Second, the five factor model has difficulty explaining the return of firms with characteristics that are different from the factor characteristics. The documented alphas combine return effects with compensation for risk dimensions that the factors do not capture, and are upper bounds on the true alignment-driven premium.

Third, the sample consists of firms that with available return data. Firms that began the sample period with high ownership but were subsequently delisted, acquired, or taken private do not appear in the data. The documented alpha is therefore conditional on survival and overstates what an investor running the strategy in real time could have earned. This survivorship bias is shared with the prior literature on ownership sorted portfolios.

Fourth, ownership is treated as a static characteristic measured at a single point in time. While the split-sample analysis indicates that this assumption is not driving the results, the approach cannot capture variation in ownership within firms over time. A dynamic specification using firm-level ownership history would be a natural extension.
